\section{The plasmonic element at terahertz frequencies}
\label{sec:plasmonic}

\subsection{Small signal: the terminations rectify, the shape tunes}
\label{sec:small}

The solver is first checked where an analytic answer exists. For a uniform
channel driven by an ideal source with an open drain, the shallow-water
closure reproduces the closed form of Dyakonov and
Shur~\cite{dyakonov1996detection} to three digits from $10\GHz$ to $3\THz$,
through the plasmon resonances at $100$~K and on the broadband plateau at
$300$~K [Fig.~\ref{fig:small}(a)]. The Dirac closure of Eq.~\eqref{eq:s2}
raises the response by $40\%$ near the resonances and by up to a factor two
between them; the difference is the nonlinearity of the Dirac fluid itself,
whose density-dependent mass and quantum capacitance add second-order terms
that the shallow-water model lacks.

Figure~\ref{fig:small}(b) shows the $1\um$ funnel. Its impedance is $80$ to
$200\ \Omega$ across the band, so it absorbs $80$ to $95\%$ of the available
power of a $75\ \Omega$ source at $100$~K. With the DS termination the
responsivity is $0.15$ to $0.34\AW$ at $100$~K and $0.07$ to $0.09\AW$ at
$300$~K, flat from $10\GHz$ to beyond $1\THz$. With the symmetric
termination, which isolates the shape, it is $10^{-5}$ to $10^{-2}\AW$ below
$100\GHz$, rises to $0.13\AW$ at $1\THz$ at $100$~K as the taper becomes
comparable to the plasmon wavelength, and stays at a few mA/W at $300$~K,
where friction and viscosity damp the plasmon. Below $100\GHz$ the
shape-only response is the curvature of the steady current-voltage relation
of Sec.~\ref{sec:dc}; a ladder test ties the two solvers together there.
The sign of the DS response follows the termination: exchanging the roles of
source and drain reverses it. That the shape of a channel modifies the
response set by its boundary conditions, rather than replacing it, is in
line with what is known for silicon and graphene field-effect detectors: a
non-uniform channel changes the responsivity by about ten per
cent~\cite{zhang2021terafet}, and a Corbino geometry, which concentrates
the field at the inner contact, raises it
substantially~\cite{khavronin2020singularity}.

The $175\nm$ funnel [Fig.~\ref{fig:small}(c)] has a dc resistance of
$1.9$~k$\Omega$ and absorbs only $15$ to $80\%$ of the available power, but
it has one feature the wide device lacks. At $1.87\THz$ the taper passes
through a series resonance, the horn is matched to the antenna, and the
shape-only responsivity reaches $0.22\AW$ at $100$~K, half of the $0.45\AW$
of the DS termination at the same frequency. The resonance belongs to the
taper and not to the leads: a control with $1\um$ leads has it at the same
frequency, with an impedance of $57-8i\ \Omega$ there. This is the one place where the
funnel shape earns its keep in a plasmonic design, as a matched horn. At
$300$~K it gives $0.03\AW$ against $0.10\AW$ for the terminations.

The Tesla valve [Fig.~\ref{fig:small}(d)] contributes almost nothing. Its
shape-only responsivity is $10^{-5}$ to $10^{-3}\AW$ below $300\GHz$ and
reaches $2\times10^{-2}\AW$ only at a $100$~K resonance near $0.7\THz$, ten
to twenty times below the funnel at every frequency and temperature. Its
resistance of $6$ to $11$~k$\Omega$ is two orders of magnitude from the
antenna, so $3$ to $50\%$ of the available power is absorbed, and whatever
the DS termination adds, $0.02$ to $0.1\AW$, is the terminations and not
the loops. The valve's asymmetry is an inertial loss that needs a Reynolds
number of ten; at terahertz frequencies the long path that produces it is
also several plasmon wavelengths long, so its stages do not act coherently.

With ideal contacts and the Johnson noise of $\Rdc$ alone, the
noise-equivalent power of the funnel elements is $4$ to $60\pWrtHz$ between
$0.3$ and $2\THz$: $17\pWrtHz$ for the $1\um$ funnel at $100$~K, about
$60\pWrtHz$ at $300$~K, and $4\pWrtHz$ for the $175\nm$ funnel at its taper
resonance at $100$~K.

\begin{figure*}
\includegraphics[width=\textwidth]{\figpath/fig4_plasmonic_small_signal.pdf}
\caption{Small-signal plasmonic response. (a) Validation on a uniform
$1\um\times10\um$ channel with an ideal source and open drain: dc drain
voltage per squared ac amplitude in units of the inverse gate overdrive. The
solver with the shallow-water closure lies on the Dyakonov-Shur closed
form~\cite{dyakonov1996detection}; the Dirac closure of
Eq.~\eqref{eq:s2} is $40\%$ higher near resonance. (b)--(d) Current
responsivity into a $75\ \Omega$ antenna for the $1\um$ funnel, the $175\nm$
funnel and the classical Tesla valve, with the drain tied to the gate
(circuit asymmetry) and with the symmetric drive that isolates the shape.
Open symbols mark a negative sign. The shaded band is $0.3$ to $1\THz$.}
\label{fig:small}
\end{figure*}

\subsection{Large signal: a square law up to the gate overdrive}
\label{sec:large}

A detector is judged by responsivity and noise, a harvester by conversion
efficiency $\eta=P_{\mathrm{dc}}/\Pav$ at the power the antenna supplies.
The framework is that of Kachorovskii
\textit{et al.}~\cite{kachorovskii2013performance} for field-effect
detectors, who define the same ratio with respect to the absorbed power,
find it quadratic in the amplitude at small signal, and show that it has an
absolute maximum of the order of ten per cent, reached below threshold where
the voltage scale is the thermal one; the change from quadratic to linear
response at an amplitude equal to the gate swing is established in theory
and experiment~\cite{gutin2012plasmonic,rudin2014theory,but2014nonlinear}.
What is specific to graphene is that a gapless monolayer has no
subthreshold regime: the only voltage scale of the gated fluid is the gate
overdrive, and the question is how large the pure-number prefactor is in a
real geometry with a real antenna.
In second-order perturbation the device is a dc Thevenin source with
open-circuit voltage $V_{\mathrm{oc}}=r\,V_{a}^{2}/V_{u}$ and resistance
$\Rdc$, where $V_{u}=m^{*}s_{0}^{2}/e$ is the natural voltage unit of the
gated fluid ($1.49$~V here, the gate overdrive up to the quantum-capacitance
correction) and $r$ is a pure number. Into a matched dc load
\begin{equation}
\eta=2r^{2}\,\frac{\Rant}{\Rdc}\,\frac{V_{a}^{2}}{V_{u}^{2}}
=\frac{16\,r^{2}\Rant^{2}}{\Rdc V_{u}^{2}}\,\Pav .
\label{eq:eta}
\end{equation}
For the $1\um$ funnel $r\approx0.1$ and $\eta/\Pav$ is $1.0$ to
$2.4$~W$^{-1}$. Since $V_{a}$ cannot usefully exceed $V_{u}$, where the
channel is depleted on the negative half cycle, Eq.~\eqref{eq:eta} already
bounds the efficiency below one per cent; what the nonlinear solver decides
is where the law fails and how. The prefactor $2r^{2}\Rant/\Rdc$ is
$0.004$ to $0.009$ at $0.1$ and $0.3\THz$.

Figure~\ref{fig:large} shows the time-domain results for the DS-terminated
$1\um$ funnel. The efficiency meets the perturbative law to $0.3\%$ at the
smallest amplitude and stays within about $10\%$ of it up to $0.9$~mW of available
power at $0.1$ and $0.3\THz$ at both temperatures. At $1\THz$ the response
saturates first, $10\%$ below the law at $150\ \mu$W at $100$~K. At
$1.8$~mW, a swing of $1.04$~V on a gate overdrive of $1.43$~V, the
efficiency is $0.14$ to $0.38\%$ and the response is $0.6$ to $0.9$ of the
square law at $100$~K; at $300$~K the overdamped device is weakly
super-quadratic, $1.10$ to $1.14$ of the law at $0.1$ and $0.3\THz$. The
density inside the device swings between $0.52$ and $1.6$ of its gated
value [Fig.~\ref{fig:large}(c)] and the solver converges at every step, so
the saturation is a property of the fluid and not a failure of the method.
Returning the dc current through the antenna instead of a matched load costs
the expected $15\%$ and nothing more.

The shape-only response behaves worse, not better, at large signal. With
the symmetric drive the funnel gives $0.56$~W$^{-1}$ at $0.3\THz$, a third
of the DS term of the same device; it leaves its square law at $150\ \mu$W
and is at $0.43$ of it at $0.9$~mW [Fig.~\ref{fig:large}(b)]. The picture
of a geometric diode that improves when driven hard is not what the fluid
does: the constriction's rectification is a resonance effect that detunes
as the density swings.

The scale of the result is set by Eq.~\eqref{eq:eta}. A beam of
$1$~W/m$^{2}$ delivers about $10$~nW to a half-wave antenna at $1\THz$,
which the element converts at $\eta\sim10^{-8}$; a full milliwatt focused
on one element is converted at $0.2\%$.

\begin{figure*}
\includegraphics[width=\textwidth]{\figpath/fig5_large_signal.pdf}
\caption{Finite-amplitude response of the DS-terminated $1\um$ funnel from
the time-domain solver. (a) Conversion efficiency into a matched dc load
against available power. (b) Efficiency relative to the square law of
Eq.~\eqref{eq:eta}, against the ac amplitude; the green curve is the
shape-only response (symmetric drive) at $0.3\THz$, $100$~K. (c) Smallest
and largest density in the device over a period, relative to the gated
density, at $100$~K.}
\label{fig:large}
\end{figure*}

\subsection{Operating point and contacts: the knee moves, the ceiling does not}
\label{sec:operating}

Equation~\eqref{eq:eta} suggests a lever. The efficiency saturates when
$V_{a}$ approaches $V_{u}$, at an available power
$P_{\mathrm{knee}}=(V_{u}/2)^{2}/8\Rant$, and $V_{u}$ is proportional to the
gate overdrive $e n_{0}/C_{g}\propto n_{0}d$. A lower density or a thinner
gate should bring the saturation down to the microwatts an antenna
collects. What it does to the peak depends on the pure number
$r^{2}\Rant/\Rdc$, which the solver has to supply. We swept
$n_{0}=5\times10^{10}$ to $10^{12}\percm$ and $d=30$, $10$, $5\nm$ with the
scattering times of Table~\ref{tab:params}, flagging points outside the
Fermi-liquid limit $\EF\geq3\kB T$ or the gating limit
$e n_{0}/C_{g}\geq4\kB T/e$; the saturation factors of Sec.~\ref{sec:large}
were carried over in units of $V_{u}$ (Appendix~\ref{app:knee}), and a second time-domain ladder at
$(5\times10^{11}\percm,10\nm)$ confirmed that they transfer, to $0.2\%$ at
$300$~K and within $14\%$ at $100$~K.

The saturation power moves as predicted [Fig.~\ref{fig:operating}(a)]: from
$0.92$~mW at the reference point to $33\ \mu$W at
$(5\times10^{11}\percm,10\nm)$ and $3\ \mu$W at $(5\times10^{10}\percm,30\nm)$,
with $\eta/\Pav$ rising from $1.9$ to $51$ and $1770$~W$^{-1}$ at $100$~K
and $0.3\THz$. The peak barely moves [Fig.~\ref{fig:operating}(b)]. Over the
valid grid it is $0.03$ to $0.9\%$ at $100$~K and at most $0.24\%$ at
$300$~K, where the reference point is already the best. A thinner gate
lowers it, because $r$ falls from $0.088$ to $0.050$ between $30$ and
$5\nm$ as the plasma speed drops and friction and quantum capacitance gain
weight; a lower density raises $r$, to $0.32$ at the edge of the
Fermi-liquid limit. A matched antenna in place of $75\ \Omega$ gives at
best $1.3\%$ at $100$~K and $0.27\%$ at $300$~K. If the mobility rather
than the scattering time is held constant as the density falls, the
low-density gain shrinks by a factor four to six. The design rule is
therefore to put the knee at the available power; the element then converts
at a few tenths of a per cent, and nothing in $n_{0}$ or $d$ raises that
number.

Contact resistance is the dominant practical loss, and it is not a series
loss. A contact resistance $R_{c}$ at each contact of the reference device
reduces $\eta/\Pav$ by factors $1.6$, $4.6$ and $78$ for $R_{c}=100$, $300$
and $1000\ \Omega\um$ at $100$~K and $0.3\THz$
[Fig.~\ref{fig:operating}(c)], the responsivity from $0.21\AW$ to $0.15$,
$0.07$ and $0.01\AW$, and it raises the NEP from $27$ to $35$, $59$ and
$240\pWrtHz$. The series picture underestimates this. The gate carries
displacement current, so the ac drain current is not the negative of the
source current, and a resistance at the drain lifts it off the gate
potential at the signal frequency: the boundary condition that makes the
Dyakonov-Shur mechanism work is partly removed. Matching the antenna to the
new impedance recovers little. That a lower access resistance raises the
photoresponse of graphene detectors strongly, by up to a factor twenty for
a tenfold reduction, has been measured and described with a
series-resistance model~\cite{castello2024impact}, and that the shunt
between gate and drain at the signal frequency is what makes a transistor
mix is the basis of the distributed self-mixing
picture~\cite{lisauskas2009rational}; the present result connects the two.

For detection the same lever is worth having
[Fig.~\ref{fig:operating}(d)]. The Johnson-limited NEP of the $1\um$ funnel
at $0.3\THz$ falls from $27\pWrtHz$ at the reference point to about
$1\pWrtHz$ at $5\times10^{10}\percm$ at $100$~K if the $5$~ps scattering
time survives there, $2\pWrtHz$ with constant mobility or with
$300\ \Omega\um$ contacts; at $300$~K the best valid point is
$15\pWrtHz$.

\begin{figure*}
\includegraphics[width=\textwidth]{\figpath/fig6_operating_point.pdf}
\caption{Carrier density, gate distance and contacts for the DS-terminated
$1\um$ funnel. (a) Saturation power $P_{\mathrm{knee}}$ against the product
$n_{0}d$; the dotted line is $(n_{0}d)^{2}$ through the reference point, and
the departure at small $n_{0}d$ is the quantum capacitance. (b) Predicted
peak efficiency at $0.3\THz$. (c) Efficiency per watt against contact
resistance, relative to ideal contacts, at the reference point. (d)
Johnson-limited NEP at $0.3\THz$. Open symbols are outside the
Fermi-liquid or gating limits.}
\label{fig:operating}
\end{figure*}
